% !TITLE 登岳阳楼
% !AUTHOR 杜甫
% !DYNASTY 唐
% !ORDER 29

\begin{poem}
昔闻洞庭水，今上岳阳楼。
吴楚东南坼\textsuperscript{1}，乾坤日夜浮\textsuperscript{2}。
亲朋无一字\textsuperscript{3}，老病有孤舟。
戎马关山北\textsuperscript{4}，凭轩涕泗流\textsuperscript{5}。
\end{poem}

\begin{annotations}
\notetext{1}{吴楚东南坼：吴地和楚地被洞庭湖从中间分裂开来。坼（chè），裂开。洞庭湖之阔，像在地理上切割开了大地。}
\notetext{2}{乾坤日夜浮：日月星辰日夜在湖水中浮动。乾坤，天地。极言湖面的浩瀚。}
\notetext{3}{无一字：没有一封信。战乱中音讯断绝，“无一字”比“无音讯”更具体。}
\notetext{4}{戎马关山北：北方关山还在打仗（当时吐蕃入侵，长安戒严）。即使身处困境，杜甫依然心系国事。}
\notetext{5}{凭轩涕泗流：倚着栏杆，眼泪鼻涕一起流下来。涕（tì），眼泪；泗（sì），鼻涕。}
\end{annotations}

\begin{appreciation}
《登岳阳楼》是杜甫晚年最重要的诗作之一。这一年是大历三年（768年），杜甫五十七岁，一身是病，坐在一只小船上。他终于登上了年轻时就想登的岳阳楼。

昔闻洞庭水，今上岳阳楼——"昔闻"和"今上"之间，隔着大半辈子。杜甫年轻时读过孟浩然写洞庭的名句："气蒸云梦泽，波撼岳阳楼。"那时他满脑子功业，大约也想过，总有一天要亲眼看看这片水。真站上来了，人已经老了。

吴楚东南坼，乾坤日夜浮——吴、楚是这一带两个古国名。洞庭湖大到什么程度？大到把吴地和楚地从中间劈开，大到日月星辰都像浮在湖面上。这是极言湖之浩瀚，气魄大得吓人。可写出这种气魄的，是一个老病缠身的老人。壮阔的天地与渺小的个人，在这一联里正面相撞。

亲朋无一字，老病有孤舟——这两句从云端一下跌回地面。无一字，不是笼统的"没有音讯"，是连一封信的一个字都没有收到，更具体，更冰凉。有孤舟，不是说他拥有这只船，是说他除了这只船，什么也没有。一个"有"字，写尽了一无所有。

戎马关山北，凭轩涕泗流——戎马，是战争。北方的关山上还在打仗，吐蕃入侵，长安戒严。凭轩，倚着栏杆。涕泗流，眼泪鼻涕一起流下来。杜甫不写"泪沾襟"这样的雅词，他写涕泗流——不美，但是真。他的身体已经这样了，心里惦记的还是北方的战事。这让我想起《离骚》里的屈原："长太息以掩涕兮，哀民生之多艰。"为苍生掉眼泪的诗人，两千年里是一脉相承的。

两年后的冬天，杜甫死在湘江的一条船上。岳阳楼是他见过的最大的风景，也是他人生最后一站。

以后你要是去湖南，替哥哥上岳阳楼看看洞庭湖。看水的时候想想杜甫：一个人可以很小，心里装的东西，可以很大。
\end{appreciation}
